\documentclass[11pt,a4paper]{article}
\usepackage[margin=1in]{geometry}
\usepackage{amsmath,amssymb}
\usepackage{booktabs}
\usepackage{hyperref}
\usepackage{xcolor}
\usepackage{float}

\definecolor{good}{rgb}{0.0,0.5,0.0}
\definecolor{moderate}{rgb}{0.7,0.5,0.0}

\title{EEG Scalp-to-In-Ear Signal Prediction:\\
  Extension 007 --- Real EEG Data (Ear-SAAD)}
\author{Automated Research Loop}
\date{April 7, 2026}

\begin{document}
\maketitle

\begin{abstract}
We evaluate scalp-to-in-ear EEG prediction on the Ear-SAAD dataset
(Geirnaert et al.\ 2025): 15 subjects with simultaneous 27-channel scalp
and 12-channel in-ear recordings at 20\,Hz (1--9\,Hz band). The closed-form
spatial filter achieves $r=0.343$ (pooled) and $r=0.326 \pm 0.092$ (LOSO).
The convolutional encoder provides negligible improvement ($r=0.348$).
These are the first real-data results for this pipeline, demonstrating that
scalp-to-in-ear prediction is feasible but substantially harder than on
synthetic data ($r=0.887$), with large inter-subject variability.
\end{abstract}

\section{Dataset}

The Ear-SAAD dataset (Zenodo record 16536441) provides simultaneous scalp,
around-ear, and in-ear EEG from 15 subjects performing an auditory attention
decoding task (competing speakers). Data is preprocessed to 1--9\,Hz at 20\,Hz
sampling rate, with artifact segments marked as NaN.

\begin{itemize}
  \item \textbf{Scalp}: 27 channels (10-20 positions, excluding mastoids M1/M2
    and reference Fp1-cr)
  \item \textbf{In-ear}: 12 channels (6 per ear: ELA, ELB, ELC, ELE, ELI, ELT,
    ERA, ERB, ERC, ERE, ERI, ERT)
  \item \textbf{Trials}: 6 per subject $\times$ 10 minutes = 60 minutes total
  \item \textbf{Windowing}: 2-second windows (40 samples) with 1-second stride
  \item \textbf{NaN handling}: Linear interpolation within channels
\end{itemize}

\section{Results}

\subsection{Pooled Evaluation}

\begin{table}[H]
\centering
\caption{Pooled results (70/15/15 chronological split across all subjects).}
\begin{tabular}{lrrr}
\toprule
Model & Mean $r$ & RMSE & SNR (dB) \\
\midrule
Closed-form & 0.343 & 0.898 & 0.52 \\
Conv Encoder & 0.348 & 0.897 & 0.54 \\
\bottomrule
\end{tabular}
\end{table}

\begin{table}[H]
\centering
\caption{Per-channel Pearson $r$ (pooled, closed-form).}
\begin{tabular}{lrlr}
\toprule
Channel & $r$ & Channel & $r$ \\
\midrule
ELA & \textcolor{good}{0.455} & ERA & 0.335 \\
ELB & 0.343 & ERB & 0.382 \\
ELC & 0.242 & ERC & 0.395 \\
ELE & 0.264 & ERE & 0.362 \\
ELI & 0.258 & ERI & \textcolor{good}{0.410} \\
ELT & 0.267 & ERT & 0.402 \\
\bottomrule
\end{tabular}
\end{table}

\subsection{Leave-One-Subject-Out (LOSO)}

\begin{table}[H]
\centering
\caption{LOSO per-subject results (closed-form).}
\begin{tabular}{rrrr|rrrr}
\toprule
Subj & $r$ & SNR & & Subj & $r$ & SNR \\
\midrule
1  & 0.293 & 0.34 &  & 9  & \textcolor{good}{0.500} & 1.08 \\
2  & \textcolor{moderate}{0.143} & $-0.15$ & & 10 & 0.268 & 0.31 \\
3  & \textcolor{good}{0.424} & 0.80 &  & 11 & 0.408 & 0.76 \\
4  & 0.240 & 0.22 &  & 12 & 0.272 & 0.36 \\
5  & 0.349 & 0.54 &  & 13 & \textcolor{good}{0.443} & 0.80 \\
6  & 0.345 & 0.55 &  & 14 & 0.265 & 0.25 \\
7  & 0.327 & 0.50 &  & 15 & 0.391 & 0.69 \\
8  & 0.230 & 0.10 &  &    &       &      \\
\bottomrule
\end{tabular}
\end{table}

\begin{table}[H]
\centering
\caption{LOSO summary: real vs.\ synthetic data.}
\begin{tabular}{lrrrr}
\toprule
& $r$ (mean) & $r$ (std) & SNR mean & SNR std \\
\midrule
Real data (LOSO) & 0.326 & 0.092 & 0.48 & 0.31 \\
Synthetic (LOSO) & 0.861 & 0.012 & 5.69 & 0.35 \\
\bottomrule
\end{tabular}
\end{table}

\section{Analysis}

\subsection{Real vs.\ Synthetic: The Gap}

The $\Delta r = 0.535$ gap between real ($r=0.326$) and synthetic ($r=0.861$)
LOSO performance reflects fundamental differences:

\begin{enumerate}
  \item \textbf{Nonlinear volume conduction}: Real skull/tissue conductivity is
    frequency-dependent and nonlinear, not a simple mixing matrix.
  \item \textbf{Subject-specific anatomy}: Each subject has unique skull geometry,
    causing inter-subject $\sigma_r = 0.092$ (vs.\ 0.012 on synthetic).
  \item \textbf{Low bandwidth}: The 1--9\,Hz preprocessed data captures only
    delta/theta activity; higher frequencies may carry more spatial information.
  \item \textbf{Independent noise}: Scalp and in-ear electrodes pick up
    different artifacts (muscle, motion), reducing correlations.
\end{enumerate}

\subsection{Conv Encoder vs.\ Closed-Form}

The conv encoder ($r=0.348$) barely outperforms the closed-form ($r=0.343$),
a $\Delta r = +0.005$ margin. This suggests:
\begin{itemize}
  \item The scalp-to-in-ear mapping is approximately linear in this frequency
    band, consistent with volume conduction physics.
  \item The 12K-parameter conv encoder cannot extract nonlinear structure
    from 2-second windows at 20\,Hz (only 40 time points).
  \item Higher sampling rates or broader frequency bands may reveal
    nonlinear dependencies.
\end{itemize}

\subsection{Per-Channel and Per-Subject Variation}

The best-predicted channels (ELA: $r=0.455$, ERI: $r=0.410$) likely have
electrode positions most proximal to scalp channels with high volume
conduction overlap. The worst (ELC: $r=0.242$) may sit deeper in the ear
canal with less scalp coupling.

Subject variability is striking: subject~9 ($r=0.500$) is 3.5$\times$
better than subject~2 ($r=0.143$). This likely reflects anatomical
differences in ear canal geometry and electrode contact quality.

\section{Significance}

\textbf{These are the first real-data results for this pipeline.} Despite
modest correlations, we have demonstrated:
\begin{enumerate}
  \item Scalp EEG \emph{does} contain information predictive of in-ear EEG
    ($r=0.34$, $p \ll 0.001$).
  \item A simple 27$\times$12 linear spatial filter captures most of the
    extractable signal.
  \item Cross-subject transfer is feasible (LOSO $r=0.326$) but variable.
  \item The synthetic pipeline was a useful development tool but
    dramatically overestimates real-world performance.
\end{enumerate}

\section{Next Steps}

\begin{enumerate}
  \item \textbf{Broadband prediction}: Obtain raw data (12.2\,GB BIDS archive)
    to test at full bandwidth (e.g., 1--45\,Hz at 256\,Hz).
  \item \textbf{Subject-specific fine-tuning}: Use a few minutes of
    target-subject data to adapt the pooled model.
  \item \textbf{Temporal context}: Use FIR filters or RNNs to capture
    cross-lag dependencies at higher sampling rates.
  \item \textbf{Around-ear channels}: Include the 19 around-ear (cEEGrid)
    channels as additional targets or features.
\end{enumerate}

\end{document}
